\chapter{System Analysis and Architecture}\label{ch:architecture}

\section{Architectural Style}

The system is a \emph{modular monolith}: one backend process and one single-page frontend, deployed together with
the database by Docker Compose (ADR-0002). The project's engineering contract forbids microservices and plugin
frameworks without a demonstrated need. A single process keeps transactions simple (a paper and all its chunks are
written in one transaction) and makes the system easy to run on one laptop.

\autoref{fig:deployment} shows the deployed components. The browser talks only to nginx, which serves the built
React application and proxies \code{/api/v1} to the FastAPI backend. The backend reaches PostgreSQL inside the
Compose network, Ollama on the host through \code{host.docker.internal}, and arXiv over HTTPS through an
allow-listed client. All published ports bind to \code{127.0.0.1}.

\begin{figure}[htbp]
  \centering
  \fitfigure{figures/arch/deployment}
  \caption{Deployment and component architecture}\label{fig:deployment}
\end{figure}

\section{Backend Layers and Responsibilities}

The backend is written in Python with FastAPI \cite{fastapi}, Pydantic schemas at every boundary, and SQLAlchemy
with Alembic migrations \cite{sqlalchemy}; the frontend uses React \cite{react} with TypeScript.
The backend is divided into packages with a strict dependency direction: \code{api} $\rightarrow$
\code{services} $\rightarrow$ \code{ingestion}, \code{retrieval}, \code{generation}, \code{evaluation}
$\rightarrow$ \code{db}. Routes never contain SQL or prompts, and the generation package never touches the
database.

\begin{table}[htbp]
  \centering
  \caption{Backend packages}\label{tab:packages}
  {\small
  \begin{tabularx}{\textwidth}{@{}lL@{}}
    \toprule
    \textbf{Package} & \textbf{Responsibility} \\
    \midrule
    \code{api/v1} & HTTP routes, Pydantic request and response schemas, error envelope \\
    \code{services} & Use cases: ingestion jobs, search, question answering, analyses, corpus, health \\
    \code{ingestion} & arXiv client, PDF validation and page-aware extraction, chunking, file storage \\
    \code{retrieval} & Embedding model wrapper, dense, lexical and hybrid search, cross-encoder re-ranker \\
    \code{generation} & Provider interface and Ollama adapter, prompts and output schemas, citation checking \\
    \code{evaluation} & Evaluation data set, metrics, run records, provenance and profiling commands \\
    \code{db} & SQLAlchemy models and sessions; schema owned by Alembic migrations \\
    \code{core} & Settings, logging, typed domain errors \\
    \bottomrule
  \end{tabularx}
  }
\end{table}

A composition root (\code{container.py}) builds the long-lived services from settings; tests build their own
containers with fakes, which keeps the services independent of their concrete dependencies.

\section{Frontend and API Boundary}

The frontend is a React and TypeScript single-page application with views for the dashboard, asking questions,
searching passages, the corpus, comparison, adding papers, history and settings. It calls the versioned REST API
through a single typed client (\code{frontend/src/api/client.ts}). Every error response uses the envelope
\code{\{"error": \{"code", "message"\}\}}, so the interface can show actionable messages. The main endpoints are
listed in \autoref{app:api}.

\section{Data Design}

Nine tables hold the persistent state (\code{docs/architecture/data-model.md}):
\begin{itemize}
  \item \code{papers}: metadata; \code{arxiv\_id} is unique, which makes arXiv import idempotent.
  \item \code{documents}: the PDF behind a paper, with its SHA-256 (unique, for duplicate detection), storage key,
        page count, and the extractor, chunker and embedding-model versions used.
  \item \code{chunks}: ordinal, \code{page\_start}, \code{page\_end}, \code{char\_start}, \code{char\_end}, token
        count, text, a 384-dimensional \code{vector} with an HNSW cosine index, and a stored \code{tsvector} with a
        GIN index for lexical search.
  \item \code{ingestion\_jobs}: kind, state, error, attempts; a partial unique index allows one in-flight job per
        source.
  \item \code{queries}, \code{answers}, \code{answer\_evidence}, \code{answer\_citations}: question history.
        Evidence rows are \emph{snapshots} of the passages given to the model, so history survives deletion of a
        paper.
  \item \code{analyses}: summaries, syntheses and extractions with their evidence snapshots.
\end{itemize}
Constraints are enforced in the database (foreign keys with cascades, unique keys, and checks on page ranges,
character spans and token counts). The schema changes only through Alembic migrations (0001 to 0006); a test
fails if the ORM models and migrations drift apart.

\section{Ingestion and Retrieval Pipelines}

\autoref{fig:dataflow} traces data from a source document to an answer citation. Ingestion validates the input,
stores the PDF under a name derived from its hash, extracts text page by page while recording the character
offset at which each page starts, chunks the text into 256-token windows with a 38-token overlap (ADR-0003),
embeds the chunks in batches and writes everything in one transaction. Retrieval embeds the question, runs dense
and lexical search, fuses the two rankings with RRF ($k=60$) and re-orders the fused candidates with the
cross-encoder (ADR-0008). The provenance chain---paper, page, character span, chunk, retrieval score, citation---is
preserved end to end.

\begin{figure}[htbp]
  \centering
  \fitfigure{figures/arch/dataflow}
  \caption{Data flow from source document to answer citation}\label{fig:dataflow}
\end{figure}

\section{Inference and Citation Flow}

Question answering keeps only passages whose cosine similarity is at least 0.30, labels them \code{P1} to
\code{Pn}, wraps them in delimiters, and sends them to Ollama together with a JSON schema that constrains the output
to claims, each listing the labels it relies on, and a status (ADR-0006). The server parses the output, resolves
every label against the set actually supplied, drops claims that consist only of labels, and returns
``insufficient evidence'' if the model abstains or if no claim has a valid citation. The detailed interaction is
shown in the sequence and communication diagrams (\autoref{fig:uml-sequence}, \autoref{fig:uml-communication}).

\section{Error Handling and External Dependencies}

Expected failures are typed domain errors mapped to HTTP status codes: invalid input 422, not found 404, conflict
409 (duplicate upload), upstream failure 502 (for example arXiv rate limiting), dependency unavailable 503 (Ollama
or the database not reachable), and generation timeout 504. Unexpected exceptions return a fixed
\code{internal\_error} message and are logged on the server only. A readiness endpoint reports the database,
embedding model, re-ranker and Ollama separately. Search keeps working when Ollama is down (NFR-16).

\section{Design Decisions}

\begin{table}[htbp]
  \centering
  \caption{Architecture decision records}\label{tab:adrs}
  {\small
  \begin{tabularx}{\textwidth}{@{}lLL@{}}
    \toprule
    \textbf{ADR} & \textbf{Decision} & \textbf{Main alternative rejected} \\
    \midrule
    0001 & Local Ollama for generation; no hosted fallback & Hosted Gemini model named in the proposal \\
    0002 & Modular monolith backend, separate SPA, monorepo & Separate ingestion worker service \\
    0003 & Chunk window of 256 tokens with 38 overlap, within the encoder input limit & Proposal's 500-token chunks, of which the encoder would see only about half \\
    0004 & Python 3.12 via uv; PostgreSQL with pgvector in Docker & Native Windows pgvector build \\
    0005 & In-process ingestion jobs backed by a job table & Celery with Redis \\
    0006 & Labelled evidence, constrained JSON, server-side citation checks & Prompt with a worked example (it produced a cited answer to an unanswerable question) \\
    0007 & Evaluation labels written by the AI assistant, with recorded provenance (team decision) & Labels written by the team (the contract's original rule) \\
    0008 & Hybrid dense and lexical retrieval with cross-encoder re-ranking & Dense retrieval only \\
    \bottomrule
  \end{tabularx}
  }
\end{table}
